\section{Basic finite energy potentials on a Hopf manifold}
\label{sec:hopf-basic-perron}

The Hopf geometry of \cref{sec:hopf-all-dimensions} gives a nonclosed background in which the potentials pulled back from the projective base carry a complete geodesic metric. The Perron envelope between two such potentials uses all subgeodesics, whereas the metric measures lengths only along paths of pulled-back potentials. The distinction matters: allowing additional smooth paths can decrease the length.

Fix $n\geq3$ and put $d=n-1$. Write
\[
 X=H_n,\qquad Y=\mathbb P^d,\qquad p=p_n,
 \qquad \omega=\omega_n=\alpha_n+\beta_n.
\]
Let $\kappa$ be the mass-one Fubini--Study form on $Y$, so $\alpha_n=p^*\kappa$. On the cover, write $r=\log|z|^2$ for the function denoted $s$ in \cref{sec:hopf-all-dimensions}; it descends to $X$ as a function defined modulo two, and $\beta_n=\mathrm dr\wedge\dc r$. The compact elliptic group
\[
 E=\mathbb C/(\mathbb Z+2\pi\mathrm i\mathbb Z)
\]
acts by scalar multiplication, preserves $\omega$, and acts transitively on each fiber. A potential is \emph{basic} if it is invariant under this action; the basic potentials are exactly the pullbacks $p^*v$ of functions $v$ on $Y$.\index{Hopf manifold!basic potentials}

These manifolds are not products of positive-dimensional compact complex manifolds. Indeed, $H_n$ is diffeomorphic to $S^{2n-1}\times S^1$, whose intermediate positive even-degree real cohomology vanishes. A product with a compact complex factor of dimension $k$, $0<k<n$, would have a nonzero degree-$2k$ class: pullback of that factor's top class is injective because projection has a smooth section. Thus the fiber reduction below takes place on a nonproduct manifold.

\Needspace{14\baselineskip}
\begin{proposition}\label{prop:hopf-basic-descent}
Let $I\subset\mathbb R$ be a nonempty open interval, let $S=S_I$ be the corresponding strip, and let $V$ be a quasi-plurisubharmonic function on $Y\times S$. Then
\[
 \operatorname{pr}_X^*\omega+\ddc(p\times\mathrm{id})^*V\geq0
 \quad\Longleftrightarrow\quad
 \operatorname{pr}_Y^*\kappa+\ddc V\geq0.
\]
For every $v\in\PSH(Y,\kappa)$, the extended non-pluripolar measure satisfies
\begin{equation}\label{eq:hopf-basic-measure}
 \MA^{\mathrm{ext}}_\omega(p^*v)
   =n\beta_n\wedge p^*\MA^{\mathrm{ext}}_\kappa(v),
\end{equation}
where $\beta_n\wedge p^*\nu$ denotes the measure $F\mapsto\int_Y\bigl(\int_{p^{-1}(y)}F\beta_n\bigr)\mathrm d\nu(y)$. In particular, for every nonnegative Borel function $F$ on $Y$,
\begin{equation}\label{eq:hopf-basic-fiber-integral}
 \int_X p^*F\,\MA^{\mathrm{ext}}_\omega(p^*v)
 =2n\int_YF\,\MA^{\mathrm{ext}}_\kappa(v).
\end{equation}
If $v_0,v_1\in\PSH(Y,\kappa)$ admit an admissible subgeodesic on $Y$, then the Perron envelope on $X$ between $p^*v_0$ and $p^*v_1$, taken over all admissible subgeodesics on $X$, is the pullback of the Perron envelope on $Y$ between $v_0$ and $v_1$.
\end{proposition}
\begin{proof}
Write $\alpha=\alpha_n$ and $\beta=\beta_n$. Locally
$\beta=(\mathrm i/(2\pi))\partial r\wedge\bar\partial r$, and $\partial r$ has a nonzero vertical component. Every vector on $Y\times S$ therefore has a lift annihilated by $\partial r$. Testing the current on these lifts removes $\beta$; integrating the resulting distributional inequality against a nonnegative compact fiber density gives positivity on the base. Conversely, pullback of a positive base current plus $\beta$ is positive. An invariant quasi-plurisubharmonic function descends canonically by restricting to local holomorphic sections, as in the proof of \cref{prop:psh-descent-submersion}, so the argument also applies to singular potentials.

For bounded $v$, the bounded products commute with pullback by \cref{lem:bounded-products-submersion}, which gives
\[
 (\omega+\ddc p^*v)^n
  =(\beta+p^*(\kappa+\ddc v))^n
  =n\beta\wedge p^*(\kappa+\ddc v)^d.
\]
All other terms vanish by $\beta^2=0$ and the dimension of $Y$. Apply this identity to $\max(v,-j)$ and restrict to $\{v>-j\}$. Increasing finite-level limits give \eqref{eq:hopf-basic-measure}, without assigning mass to the pole set. The identity $\int_{p^{-1}(y)}\beta=2$ from \cref{sec:hopf-all-dimensions} gives \eqref{eq:hopf-basic-fiber-integral}.

The family of all admissible subgeodesics on $X$ between $p^*v_0$ and $p^*v_1$ is preserved by $E$, so its supremum and the upper semicontinuous regularization are invariant. A pulled-back admissible subgeodesic of $Y$ belongs to this family, so the envelope is the largest admissible subgeodesic by \cref{prop:hermitian-perron}. By the first assertion, this invariant envelope descends to an admissible subgeodesic on $Y$, since its joint upper limits at the endpoints descend as well. Conversely, every admissible subgeodesic on $Y$ pulls back to one on $X$. The two envelopes consequently agree; in submersion charts, upper semicontinuous regularization commutes with pullback.
\end{proof}

Let $\mathcal E^2(Y,\kappa)$ denote the set of $v\in\PSH(Y,\kappa)$ of non-pluripolar mass one with finite squared energy $\int_Y|v|^2\MA^{\mathrm{ext}}_\kappa(v)$. Define
\[
 \mathcal B_\infty=\{p^*v:v\in C^\infty(Y),\ \kappa+\ddc v>0\},
 \qquad
 \mathcal B_2=p^*\mathcal E^2(Y,\kappa).
\]
Thus $\mathcal B_2$ consists precisely of the basic potentials of extended mass $2n$ with $\int_X|u|^2\MA^{\mathrm{ext}}_\omega(u)<\infty$. On $\mathcal B_\infty$, let $d_B$ be the infimum of the lengths of smooth paths in $\mathcal B_\infty$ for the $L^2$ metric $G$ of the previous section.\index{Finite energy!basic Hopf potentials}

A geodesic metric space $(Z,\delta)$ is CAT$(0)$ if every constant-speed geodesic $\gamma\colon[0,1]\to Z$ satisfies
\[
 \begin{aligned}
 \delta(z,\gamma_t)^2\leq{}&
 (1-t)\delta(z,\gamma_0)^2+t\delta(z,\gamma_1)^2\\
 &-t(1-t)\delta(\gamma_0,\gamma_1)^2
 \end{aligned}
\]
for $z\in Z$ and $0\leq t\leq1$. In particular, its geodesics are unique.

We use the classical K\"ahler completion theorem: the length metric of the $L^2$ metric $\int_Y\xi_1\xi_2(\kappa+\ddc v)^d$ has completion $\mathcal E^2(Y,\kappa)$, the completed metric is CAT$(0)$, and its geodesic between two points is their Perron segment. That segment attains its endpoints in capacity, and all its slices lie in $\mathcal E^2$. We also use the full-mass characterization $P_\kappa[v]=0$ for $v\in\mathcal E^2$. These are closed-background inputs on $Y$; their exact source and operator normalization are specified in \cref{history:hopf-basic-perron}.

\Needspace{12\baselineskip}
\begin{theorem}\label{thm:hopf-basic-completion}\index{Metric completion!basic Hopf potentials}
The completion of $(\mathcal B_\infty,d_B)$ is $\mathcal B_2$. It is a CAT$(0)$ space, and its unique constant-speed geodesic between two points is their full Hermitian Perron segment. This segment attains both endpoints in $L^1(X,\mathrm dV)$ for every smooth positive volume form $\mathrm dV$. For every $u\in\mathcal B_2$ and $c\in\mathbb R$,
\[
 d_B(u,u+c)=\sqrt{2n}|c|.
\]
\end{theorem}
\begin{proof}
Fiber integration gives the exact restricted metric
\begin{equation}\label{eq:hopf-basic-metric}
 G_{p^*v}(p^*\xi_1,p^*\xi_2)
 =2n\int_Y\xi_1\xi_2(\kappa+\ddc v)^d.
\end{equation}
The K\"ahler completion theorem therefore gives the stated completion and CAT$(0)$ property after multiplying base distances by $\sqrt{2n}$. If $\widetilde d_D$ denotes the completed metric written with operator $\mathrm i\partial\bar\partial$ and background $\kappa$, the normalization is explicitly
\begin{equation}\label{eq:hopf-basic-distance}
 d_B(p^*v,p^*w)
 =2\pi\sqrt{2n}\,\widetilde d_D(v/(2\pi),w/(2\pi)),
\end{equation}
because $\kappa+\mathrm i\partial\bar\partial(v/(2\pi))=\kappa+\ddc v$.
The geodesic between two points of $\mathcal B_2$ is the pullback of their Perron segment on $Y$, and \cref{prop:hopf-basic-descent} identifies it with their Perron segment on $X$, taken over all admissible subgeodesics. The base segment attains its endpoints in capacity, hence in $L^1(Y)$. Since
\[
 \int_X|p^*(v-w)|\frac{\omega^n}{2n}
 =\int_Y|v-w|\kappa^d,
\]
the same holds on $X$; every smooth positive $\mathrm dV$ has bounded density ratio with $\omega^n$ on compact $X$. For smooth $v$, the constant translation is a base geodesic with speed $|c|$. Smooth metric approximation and translation by the same constant extend the distance formula to every $v\in\mathcal E^2$.
\end{proof}

For an absolutely continuous curve $\gamma\colon[0,1]\to\mathcal B_2$, define its metric action by
\[
 \mathcal A_B(\gamma)=\frac12\int_0^1|\dot\gamma|_{d_B}^{\,2}\,\mathrm dt.
\]
Here the metric derivative is the almost everywhere defined limit
\[
 |\dot\gamma_t|_{d_B}
 =\lim_{h\to0}\frac{d_B(\gamma_{t+h},\gamma_t)}{|h|}.
\]
The action is defined even when the potentials are unbounded; it does not require a Riemannian tensor at a singular potential. To distinguish them from $\mathcal A_B$, we call the action $\mathscr S$ of a smooth path in $\HH_\omega$, defined in the previous section, its \emph{raw action}, and its length for $G$ its \emph{raw length}.

\Needspace{12\baselineskip}
\begin{proposition}\label{prop:hopf-basic-action}\index{Action!basic Hopf potentials}
The unique action minimizer between two points of $\mathcal B_2$ is their constant-speed Perron segment, and its action is half their squared distance. On every smooth path in $\mathcal B_\infty$, $\mathcal A_B$ equals the raw action $\mathscr S$. Every smooth basic Perron segment with strictly positive spatial forms is stationary for $\mathscr S$ under all smooth fixed-endpoint variations on $X$.
\end{proposition}
\begin{proof}
The length inequality and Cauchy--Schwarz give
\[
 d_B(\gamma_0,\gamma_1)^2
 \leq\left(\int_0^1|\dot\gamma|_{d_B}\,\mathrm dt\right)^2
 \leq2\mathcal A_B(\gamma).
\]
Equality holds for the constant-speed geodesic. Equality for a connecting curve forces constant speed and equality between length and endpoint distance, hence forces a geodesic; uniqueness follows from the CAT$(0)$ property.

For a smooth basic path $p^*v_s$, put $h=v_{s+\delta}-v_s$. On the base Perron segment from $v_s$ to $v_{s+\delta}$, parametrized by $0\leq\tau\leq1$, the chord $v_s+\tau h$ bounds above, and
\[
 v_s+\tau h-C\delta^2\tau(1-\tau)
\]
bounds below. Indeed, the spatial forms along the chord have a uniform positive lower bound, $\mathrm dh=O(\delta)$, and \cref{prop:hermitian-schur-geodesic} makes the displayed path a subgeodesic for a fixed sufficiently large $C$. The right initial velocity of the Perron segment is therefore $h+O(\delta^2)$ uniformly. The classical smooth-endpoint K\"ahler distance formula gives
\[
 \lim_{\delta\to0}\frac{d_B(p^*v_s,p^*v_{s+\delta})^2}{\delta^2}
 =2n\int_Y\dot v_s^{\,2}(\kappa+\ddc v_s)^d.
\]
Integration proves $\mathcal A_B=\mathscr S$ along this path.

For stationarity, let $u_t=p^*v_t$ be a smooth basic Perron segment with strictly positive spatial forms; by \cref{prop:hopf-basic-descent}, $v_t$ is the Perron segment on $Y$ between its endpoints. Let $L(w)$ be the derivative of $\mathscr S(u+\varepsilon w)$ at $\varepsilon=0$. Uniform spatial positivity allows every smooth endpoint-zero variation $w$. Group invariance of the action and of the path gives $L(g^*w)=L(w)$ for $g\in E$. Average over Haar probability measure to obtain
\[
 L(w)=L\left(\int_Eg^*w\,\mathrm dg\right)=L(p^*q)
\]
for a smooth endpoint-zero base variation $q$. Fiber integration reduces this to $2n$ times the closed-base first variation. With $\eta_t=\kappa+\ddc v_t$, integration by parts gives that variation as
\[
 \int_0^1\int_Yq_t\left[-\ddot v_t\eta_t^d
 +d\,\mathrm d\dot v_t\wedge\dc\dot v_t\wedge\eta_t^{d-1}\right]\mathrm dt.
\]
The bracket vanishes because the base segment solves the homogeneous equation \eqref{eq:hermitian-hcma-smooth} on $Y$. Indeed, if the $(d+1)$st power of its semipositive product form were positive at some point, that form would be strictly positive near the point, and adding a small nonzero nonnegative bump function of $(y,\Real z)$ supported near the point would give a larger admissible subgeodesic. Thus $L(w)=0$ for every variation on $X$.
\end{proof}

The metric action is minimized only among curves in $\mathcal B_2$. Stationarity holds under all smooth variations on $X$, but it does not imply minimization among all smooth Hermitian paths, and \cref{sec:hopf-action-saddles} shows that minimization fails. Likewise, the initial-velocity formula used above requires smooth endpoints; it is not a definition of the metric on all of $\mathcal B_2$.

\section{Negative directions of the raw action}
\label{sec:hopf-action-saddles}

On this Hopf background, translations by constants are stationary Perron segments, and radial variations show that they do not minimize the raw action. Every smooth two-periodic function of $r$ defines a smooth function on $X$, and for a radial path $u(t,r)$,
\[
 \omega+\ddc u=(1+u_r)\alpha_n+(1+u_{rr})\beta_n.
\]
For $\langle F\rangle=\frac12\int_0^2F(r)\,\mathrm dr$, fiber integration therefore gives
\begin{equation}\label{eq:hopf-radial-action}
 \mathscr S(u)=n\int_0^1
 \left\langle u_t^2(1+u_r)^d(1+u_{rr})\right\rangle\mathrm dt.
\end{equation}
Changing the local base section shifts the start of the full period and leaves the average unchanged.

\Needspace{12\baselineskip}
\begin{theorem}\label{thm:hopf-action-saddle}\index{Action!Hopf saddle}
For every nonzero constant $c$, the full Perron segment $u_t=ct$, $0\leq t\leq1$, is stationary for the raw action, and the Hessian of the raw action at this segment is negative definite on an infinite-dimensional space of smooth fixed-endpoint variations. The infimum of raw lengths over all smooth strictly positive paths from $0$ to $c$ is strictly smaller than $d_B(0,c)=\sqrt{2n}|c|$.
\end{theorem}
\begin{proof}
The path is basic, equals the chord, and is the full Perron segment by \cref{cor:translations-gauduchon}. Since $\omega$ is Gauduchon, the same corollary shows that it is stationary; this also follows from \cref{prop:hopf-basic-action}.
Let $w(t,r)$ be smooth, two-periodic in $r$, and zero at $t=0,1$. The paths $ct+\varepsilon w$ have positive spatial forms for sufficiently small $|\varepsilon|$. Differentiating \eqref{eq:hopf-radial-action} twice in $\varepsilon$ at $\varepsilon=0$ gives
\[
 2n\int_0^1\left\langle
 w_t^2+2cw_t(dw_r+w_{rr})
 +c^2\left(\frac{d(d-1)}2w_r^2+dw_rw_{rr}\right)
 \right\rangle\mathrm dt.
\]
The integral of $w_rw_{rr}$ over a period vanishes, while
$\int_0^1\langle w_tw_{rr}\rangle\mathrm dt
=-\frac12[\langle w_r^2\rangle]_0^1=0$.
Completing the square gives the Hessian
\begin{equation}\label{eq:hopf-action-hessian}
 2n\int_0^1\left\langle
 (w_t+dcw_r)^2-\frac{d(d+1)}2c^2w_r^2
 \right\rangle\mathrm dt.
\end{equation}
For a positive integer $k$, take
\[
 w_k(t,r)=\sin(\pi t)\cos\bigl(\pi k(r-dct)\bigr).
\]
Then $w_{k,t}+dcw_{k,r}=\pi\cos(\pi t)\cos(\pi k(r-dct))$, and \eqref{eq:hopf-action-hessian} is
\[
 \frac{n\pi^2}2\left(1-\frac{d(d+1)}2c^2k^2\right).
\]
It is negative for all sufficiently large $k$. For distinct $k$, the functions $w_{k,t}+dcw_{k,r}$ are orthogonal under the $r$ average, and so are the derivatives $w_{k,r}$. Hence \eqref{eq:hopf-action-hessian} is negative on every nonzero finite combination $w$ of the modes $w_k$ with $k$ this large. For sufficiently small $\varepsilon\neq0$, the path $ct+\varepsilon w$ is strictly positive and has raw action strictly below $\mathscr S(ct)=nc^2$. By the Cauchy--Schwarz inequality on the unit interval, its raw length $\ell$ satisfies $\ell^2\leq2\mathscr S(ct+\varepsilon w)<2nc^2$, so $\ell<\sqrt{2n}|c|$.
\end{proof}

These competing paths need not be subgeodesics. The strict length inequality does not establish definiteness or completeness of the unrestricted length infimum. For $n=3$ and $c=1$, the mode $k=1$ gives the concrete path
\[
 u_\varepsilon(t,r)=t+\varepsilon\sin(\pi t)\cos\bigl(\pi(r-2t)\bigr).
\]
Its spatial coefficients are at least $1-\pi|\varepsilon|$ and $1-\pi^2|\varepsilon|$. Its exact action is
\begin{equation}\label{eq:hopf-saddle-polynomial}
 \mathscr S(u_\varepsilon)
 =3-\frac{3\pi^2}2\varepsilon^2+\frac{111\pi^4}{64}\varepsilon^4.
\end{equation}
Indeed, put $f=\sin(\pi t)$, $g=f'$, and $\vartheta=\pi(r-2t)$. Substitute
$u_t=1+\varepsilon(g\cos\vartheta+2\pi f\sin\vartheta)$,
$u_r=-\varepsilon\pi f\sin\vartheta$, and $u_{rr}=-\varepsilon\pi^2f\cos\vartheta$ into \eqref{eq:hopf-radial-action}. Odd powers have zero periodic average. The averaged quadratic coefficient is $(g^2-3\pi^2f^2-2\pi^2fg)/2$, and the quartic coefficient is $\pi^2f^2(g^2+6\pi^2fg+12\pi^2f^2)/8$. The identities
\[
 \begin{gathered}
 \int_0^1fg\,\mathrm dt=\int_0^1f^3g\,\mathrm dt=0,
 \\
 \int_0^1f^2\,\mathrm dt=\frac12,
 \quad \int_0^1g^2\,\mathrm dt=\frac{\pi^2}2,\\
 \int_0^1f^4\,\mathrm dt=\frac38,
 \quad \int_0^1f^2g^2\,\mathrm dt=\frac{\pi^2}8
 \end{gathered}
\]
give \eqref{eq:hopf-saddle-polynomial}. Taking $\varepsilon=1/100$ gives a positive path with action $2.9985212487\ldots<3$.

\Needspace{12\baselineskip}
\section{Explicit basic geodesics and singular potentials}

Diagonal projective automorphisms give smooth nonconstant geodesics. For real $a_0,\ldots,a_{n-1}$, set
\[
 v_t([z])=\log\frac{\sum_{j=0}^{n-1}\mathrm e^{ta_j}|z_j|^2}
 {\sum_{j=0}^{n-1}|z_j|^2}.
\]
The family is spatially nonconstant when the $a_j$ are not all equal.

\Needspace{10\baselineskip}
\begin{proposition}\label{prop:hopf-diagonal-geodesic}\index{Geodesic!diagonal projective family}
The pullback $p^*v_t$, $0\leq t\leq1$, is the full Hermitian Perron segment from $0$ to $p^*v_1$, and it is stationary under all smooth fixed-endpoint variations. The distance between its endpoints and its action are
\[
 \begin{aligned}
 d_B(0,p^*v_1)^2&=\frac{2[(\sum_ja_j)^2+\sum_ja_j^2]}{n+1},\\
 \mathscr S(p^*v)&=\mathcal A_B(p^*v)
 =\frac{(\sum_ja_j)^2+\sum_ja_j^2}{n+1}.
 \end{aligned}
\]
\end{proposition}
\begin{proof}
Let $s$ be the coordinate on the strip $S_{(0,1)}$, and let $V([z],s)=v_{\Real s}([z])$ be the complexification. Since $|\mathrm e^{a_js/2}|^2=\mathrm e^{a_j\Real s}$, the holomorphic map
$F([z],s)=[\mathrm e^{a_0s/2}z_0:\cdots:\mathrm e^{a_{n-1}s/2}z_{n-1}]$
satisfies $\operatorname{pr}_Y^*\kappa+\ddc V=F^*\kappa$. The form $F^*\kappa$ has rank $d$, strictly positive spatial slices, and zero $(d+1)$st power. The classical uniqueness for the bounded K\"ahler Dirichlet problem identifies $v_t$ with the base Perron segment from $0$ to $v_1$; \cref{prop:hopf-basic-descent,prop:hopf-basic-action} give the full segment and its stationarity.

For $x_j=|z_j|^2/\sum|z_j|^2$, the initial velocity is $\sum_ja_jx_j$. A standard complex Gaussian vector has independent exponential squared coordinates, and its projective direction has distribution $\kappa^d$. Dividing those coordinates by their sum gives the uniform simplex distribution, with moments
\[
 \int_Yx_i^2\kappa^d=\frac2{n(n+1)},
 \qquad \int_Yx_ix_j\kappa^d=\frac1{n(n+1)}\quad(i\ne j).
\]
The smooth distance formula and \eqref{eq:hopf-basic-metric} give the claimed squared distance; the action is half of it.
\end{proof}

For instance, $n=3$ and $(a_0,a_1,a_2)=(1,-1,0)$ give squared distance one and action $1/2$.

The completed space also contains potentials with different ordinary singularity types. Put $o=[1:0:\cdots:0]\in Y$ and
\[
 \psi=\log\frac{\sum_{j=1}^d|z_j|^2}{\sum_{j=0}^d|z_j|^2},
 \qquad v_a=1-(1-\psi)^a,\qquad0<a<1,
\]
with the canonical value $v_a(o)=-\infty$.

\Needspace{12\baselineskip}
\begin{proposition}\label{prop:hopf-singular-basic-energy}\index{Finite energy!radial exponent threshold}
Every $v_a$ has non-pluripolar mass one, and
\[
 v_a\in\mathcal E^2(Y,\kappa)
 \quad\Longleftrightarrow\quad a<\frac d{d+2}.
\]
At equality its squared energy diverges logarithmically. The pullback has mass $2n$ and pole set $p^{-1}(o)$. If $0<a<b<d/(d+2)$, the two pullbacks have different ordinary singularity types and finite positive $d_B$ distance.
\end{proposition}
\begin{proof}
In the chart $z_0=1$, write $s=\log|z'|^2$, $L=\log(1+\mathrm e^s)$, $P=\mathrm e^s/(1+\mathrm e^s)$, and $T=1+L-s\geq1$. The local potential $\Phi=L+1-T^a$ has derivative
\[
 A=P+a(1-P)T^{a-1},\qquad
 A'=P'(1-aT^{a-1})+a(1-a)(1-P)^2T^{a-2}>0.
\]
Both $A$ and $A'$ are positive, so the radial complex Hessian is positive. Its bounded-above local potential extends plurisubharmonically across $o$ by \cref{lem:qpsh-removable-analytic}. The function extends smoothly across the hyperplane at infinity. Its bounded plurisubharmonic truncations decrease to a nontrivial canonical potential: near $o$, $|v_a|=O((1-s)^a)$ is integrable against radial smooth volume $O(\mathrm e^{ds})\,\mathrm ds$.

The radial mass flux is $Q(s)=A(s)^d$. Its sphere coefficient is one, since $\dc\log|z'|^2$ integrates to one on the Hopf circle and $\ddc\log|z'|^2$ induces mass-one Fubini--Study on $\mathbb P^{d-1}$. Stokes on an annulus gives mass $Q(s_2)-Q(s_1)$. Since $Q(-\infty)=0$ and $Q(+\infty)=1$, finite-level locality and annular exhaustion give non-pluripolar mass one. The pole has no atom, and smoothness gives no mass on the hyperplane at infinity.

As $s\to-\infty$,
\[
 |v_a|^2\sim T^{2a},\qquad
 Q'(s)\sim d a^d(1-a)T^{-d(1-a)-1},\qquad T\sim1-s.
\]
The squared energy therefore converges exactly when $2a-d(1-a)<0$. At equality the integrand has order $T^{-1}$, giving logarithmic divergence. All other compact regions have smooth potentials and measures. Fiber integration gives the pullback assertions. Finally, $v_a-v_b=T^b-T^a\to+\infty$ at $o$, so mutual bounded comparison fails. Both points lie in the completed metric space and are distinct, giving finite positive distance.
\end{proof}

\Needspace{16\baselineskip}
\section{Model projection and test curves}
\label{sec:hopf-basic-models}

The distance just constructed retains information about the potentials that their ordinary model projection forgets. This distinction connects the example with the singularity spaces studied earlier and with the Legendre construction in the next chapter.

\Needspace{12\baselineskip}
\begin{proposition}\label{prop:hopf-basic-model-projection}\index{Model envelope!basic Hopf potentials}
For every $v\in\PSH(Y,\kappa)$,
\[
 P_\omega[p^*v]=p^*P_\kappa[v].
\]
Consequently $P_\omega[u]=0$ for every $u\in\mathcal B_2$. Pullback commutes with the forward and backward Legendre transforms \eqref{eq:hermitian-legendre-forward} and \eqref{eq:hermitian-legendre-backward} of envelope test curves on $Y$ (\cref{def:hermitian-envelope-testcurve}). Let $(v_t)_{t\geq0}$ be a family in $\PSH(Y,\kappa)$ whose restriction to every finite interval is a Perron segment. If every $v_t$, including $v_0$, lies in $\mathcal E^2(Y,\kappa)$, then each finite segment of its pullback is a $d_B$ geodesic and minimizes the metric action with its prescribed endpoints.
\end{proposition}
\begin{proof}
For finite $C$, the obstacle $p^*\min(v+C,0)$ and the family of all competitors for its envelope on $X$ are invariant under $E$. The same argument as in \cref{prop:hopf-basic-descent} gives
\[
 P_\omega(p^*\min(v+C,0))
 =p^*P_\kappa(\min(v+C,0)).
\]
Increasing limits and their upper semicontinuous regularizations commute with pullback in submersion charts, proving the model identity. The classical full-mass characterization on $Y$ gives $P_\kappa[v]=0$ for every $v\in\mathcal E^2$, hence the assertion for $\mathcal B_2$.

Pullback preserves the concavity and monotonicity in the parameter and the cutoff of an envelope test curve on $Y$, and, by the model identity, the envelope condition and the initial potential. It commutes with pointwise infima, and the submersion property makes it commute with upper semicontinuous regularization of suprema. Thus it commutes with both Legendre transforms. Finally, \cref{prop:hopf-basic-descent} identifies the pullback of every finite segment of $(v_t)$ with the full Hermitian Perron segment; its finite-energy endpoints and \cref{thm:hopf-basic-completion,prop:hopf-basic-action} give the geodesic and minimizing assertions.
\end{proof}

For endpoints of the same ordinary singularity type, \cref{thm:ordinary-type-perron} already gives their common type throughout the segment. The mass formula \eqref{eq:hopf-basic-fiber-integral} computes the mass of a chosen basic potential; it does not compute the lower or upper volume of its entire ordinary type, which also includes nonbasic potentials. A separating complete metric on all ordinary models, with compatibility under modifications and changes of background, remains open.
